\section{The rank and torsion of the chosen curve}\label{app:rank-one}

\begin{proof}[Proof of Proposition~\ref{prop:rank-one}]
For \(E:y^2=x^3-25x\), we first prove \([E(\QQ):2E(\QQ)]\le8\), which will force
rank at most one.  We then exhibit a nontorsion point and determine
the full torsion subgroup.  For the rank bound we use the standard
two-isogeny descent \cite[Chapter~III, Example~4.5, and
Chapter~X, Section~4, Example~4.8 and Proposition~4.9]{silverman}, giving the kernel and index
calculations explicitly.  For rational
\(a,b\) with \(b(a^2-4b)\ne0\), write
\[
 E_{a,b}:y^2=x^3+ax^2+bx,\qquad
 b'=a^2-4b,\qquad E'=E_{-2a,b'}.
\]
Both curves are nonsingular.  The following maps, with the exceptional
points \(O\) and \((0,0)\) sent to \(O\), are isogenies:
\begin{align}
 \phi:E_{a,b}&\longrightarrow E',&
 \phi(x,y)&=\left(\frac{y^2}{x^2},
                  y\left(1-\frac b{x^2}\right)\right),
 \label{eq:phi-descent}\\
 \psi:E'&\longrightarrow E_{a,b},&
 \psi(X,Y)&=\left(\frac{Y^2}{4X^2},
                  \frac Y8\left(1-\frac{b'}{X^2}\right)\right).
 \label{eq:psi-descent}
\end{align}
Here are the algebraic and geometric checks needed for this assertion.
For the first map, put \(X=x+a+b/x\).  Substitution gives
\[
 X\bigl((X-a)^2-4b\bigr)
     =y^2\left(1-\frac b{x^2}\right)^2,
\]
which is the target equation.  Applying the same formula to \(E'\)
gives \(E_{4a,16b}\).  The scaling \((x,y)\mapsto(4x,8y)\)
is an isomorphism from \(E_{a,b}\) to \(E_{4a,16b}\);
composing with its inverse gives
\eqref{eq:psi-descent}.

These rational formulas extend over the indicated exceptional points.
At \((0,0)\), the equation
\(y^2=x(b+ax+x^2)\) shows that \(y\) has order one and \(x\)
has order two.  The target coordinates in \eqref{eq:phi-descent}
therefore have poles of orders two and three.  Their ratios \(X/Y\)
and \(1/Y\), which are coordinates in the target chart at infinity,
are regular and vanish there.  At \(O\), the functions \(x,y\)
have poles of orders two and three, and the same argument applies.
Away from these points the formulas are regular.  The argument also
applies to the map from \(E'\), and scaling preserves regularity.
Thus the two maps are nonconstant morphisms preserving the identity;
they are group homomorphisms by the standard isogeny theorem
\cite[Chapter~III, Theorem~4.8]{silverman}.

Write \([c]\) for the class of \(c\in\QQ^*\) in
\(\QQ^*/\QQ^{*2}\).  Define
\[
 \alpha:E_{a,b}(\QQ)\longrightarrow\QQ^*/\QQ^{*2},\qquad
 \alpha(O)=[1],\quad \alpha((0,0))=[b],\quad
 \alpha((x,y))=[x]\quad(x\ne0).
\]
This is a homomorphism.  Indeed, a nonvertical line
\(y=\lambda x+\nu\) meets the cubic, with multiplicity, at roots of
\[
 x^3+(a-\lambda^2)x^2+(b-2\lambda\nu)x-\nu^2.
\]
If \(\nu\ne0\), the product of the three roots is \(\nu^2\).
If \(\nu=0\), the root at \((0,0)\) is simple, since \(b\ne0\),
and the other two roots have product \(b\).  In either case the
product of the three \(\alpha\)-values is trivial.  The third
intersection is the negative of the sum of the first two, and
negation leaves \(\alpha\) unchanged.  Vertical pairs have product
class one; pairs involving \(O\) are immediate.  This also treats
tangencies, with intersection multiplicity understood.

We next prove the exact kernel identity
\begin{equation}\label{eq:descent-kernel}
 \ker\alpha=\psi E'(\QQ).
\end{equation}
An affine image other than \((0,0)\) has square x-coordinate by
\eqref{eq:psi-descent}.  An image equal to \((0,0)\) must have
\(X\ne0\), \(Y=0\), and hence
\(X=a\pm2\sqrt b\).  Such a rational preimage exists exactly
when \(b\) is a rational square; its x-coordinate is nonzero since
\(b'\ne0\).  Thus every image is in the kernel, including the
exceptional points.

Conversely, suppose \(x=s^2\ne0\) for a point \((x,y)\) in
the kernel.  Set
\[
 X=2x+a+2y/s,\qquad Y=2sX.
\]
The product of \(2x+a+2y/s\) and \(2x+a-2y/s\) is \(b'\),
so \(X\ne0\).  The identity
\(X^2-(2a+4x)X+b'=0\) implies
\(Y^2=X^3-2aX^2+b'X\).  Formula~\eqref{eq:psi-descent}
then gives first coordinate \(x\) and second coordinate
\(s(X-b'/X)/4=y\).  When \((0,0)\) is in the kernel, its
preimages were just exhibited; \(O\) is always an image.
This proves \eqref{eq:descent-kernel}.

Let \(\alpha'\) be the analogous squareclass homomorphism on
\(E'\).  Applying the kernel calculation to \(E'\), and using
the scaling above, gives
\(\ker\alpha'=\phi E_{a,b}(\QQ)\).
For \(xy\ne0\), the x-coordinate of \(\psi\phi(x,y)\) is
\[
 \frac{(x^2-b)^2}{4y^2}=x([2](x,y)),
\]
as follows by expanding the usual tangent formula
\(((3x^2+2ax+b)/(2y))^2-a-2x\).
The points therefore agree up to sign.  At \(O\) and the points
with \(y=0\), the composition and doubling both give \(O\).
Since the image of a homomorphism is closed under negation, this
pointwise sign ambiguity still gives
\(\psi\phi E_{a,b}(\QQ)=2E_{a,b}(\QQ)\).
The map induced by \(\psi\) from \(E'(\QQ)/\phi E_{a,b}(\QQ)\)
onto \(\psi E'(\QQ)/2E_{a,b}(\QQ)\) consequently yields
\begin{equation}\label{eq:descent-index}
 [E_{a,b}(\QQ):2E_{a,b}(\QQ)]
 \le |\operatorname{im}\alpha|\,
      |\operatorname{im}\alpha'|.
\end{equation}

The descent has reduced the desired index bound to two squareclass
computations. We now bound those images for our concrete curve;
finite generation will then turn the index bound into a rank bound.
We specialize to \(a=0,b=-25\).  For integral \(a,b\), a
point with \(xy\ne0\) can have an odd valuation in its
x-coordinate only at a prime dividing \(b\).  Write \(v_p(x)\)
for the exponent of the prime \(p\) in the rational number \(x\).
To see the assertion, if
\(v_p(x)<0\), then \(x^2\) is the unique term of smallest
valuation in \(x^2+ax+b\), so
\(2v_p(y)=3v_p(x)\).  If \(v_p(x)>0\) and \(p\nmid b\),
then \(x^2+ax+b\) is a unit, so \(2v_p(y)=v_p(x)\).
In either situation \(v_p(x)\) is even.  Including the
exceptional points, this gives
\[
 \operatorname{im}\alpha\subseteq\{[1],[-1],[5],[-5]\}.
\]
The second curve is \(E':y^2=x^3+100x\).  Every nonzero
x-coordinate on it is positive, and the same valuation argument gives
\[
 \operatorname{im}\alpha'\subseteq\{[1],[2],[5],[10]\}.
\]
Its image does not contain \([2]\).  Otherwise write
\(x=2m^2/e^2\), with nonzero coprime integers \(m,e\).
Substitution gives
\[
 \left(\frac{ye^3}{2m}\right)^2=2m^4+50e^4=N^2
\]
for an integer \(N\), because a rational number with integral
square is integral.  Reduction modulo five first forces
\(5\mid m,N\).  Writing \(m=5m_1\), \(N=5N_1\) then gives
\[
 N_1^2=50m_1^4+2e^4,
\]
which is again impossible modulo five, since \(5\nmid e\).
A subgroup of the four displayed squareclasses that omits \([2]\)
has size at most two.  Thus \eqref{eq:descent-index} gives
\([E(\QQ):2E(\QQ)]\le8\).

By the Mordell--Weil theorem, \(E(\QQ)\) is finitely generated
\cite[Chapter~VIII, Theorem~4.1]{silverman}.  If its rank is \(r\)
and its finite torsion subgroup is \(\mathcal T\), then
\[
 [E(\QQ):2E(\QQ)]
   =2^r|\mathcal T/2\mathcal T|
   =2^r|\mathcal T[2]|=2^{r+2},
\]
since its rational points of order dividing two are precisely
\(O,(0,0),(5,0),(-5,0)\).  Hence \(r\le1\).
The point \((-4,6)\) is nontorsion: the Nagell--Lutz theorem
for an integral short Weierstrass equation \(y^2=x^3+Ax+B\)
says that a nonidentity
rational torsion point has integral coordinates and either
\([2]Q=O\) or \(y(Q)^2\mid4A^3+27B^2\)
\cite[Chapter~VIII, Corollary~7.2]{silverman}.
Here \(36\nmid62500\), so the rank is one.

The same theorem determines the torsion completely.  For integral
coordinates on our curve, reduction modulo three gives
\(y^2\equiv x^3-x\equiv0\), so \(3\mid y\).
A torsion point with \(y\ne0\) would also satisfy
\(y^2\mid62500\), a contradiction.  Thus every affine torsion
point has \(y=0\), giving the claimed list.  Finally choose a
point \(P\) whose image generates the infinite cyclic quotient
\(E(\QQ)/\mathcal T\).  Its cosets, and then the torsion
translates within each coset, give the asserted unique expression.
\end{proof}
